\%============================================================
\chapter{Bai Hoc: Khoan Dung --- Suc Manh Cua Su Tha Thu (Forgiveness --- The Greatest Power)}
\begin{center}
  \textit{\color{truyengold}``Forgiveness is not weakness; it is the courage of the strong.''}\\[4pt]
  \textit{(Tha thu khong phai la yeu duoi; do la long can dam cua nguoi manh.)}\\[4pt]
  \small\color{truyenblue}--- Co Kim Dong Tay ---
\end{center}
\ngancach

\section{Lincoln Va Ke Thu Chinh Tri}
\begin{truyen}{Lincoln Va Ke Thu Chinh Tri}{Lincoln --- Hoa Ky 1860}
\chuhoa{A}{}A braham Lincoln sau khi dac cu Tong Thong da lam dieu ma ca nuoc My khong ai ngo toi.\\
\textit{(Abraham Lincoln, after winning the presidency, did something no one in America expected.)}

Ong bo nhiem Edwin Stanton --- nguoi da cong khai che gieu ong la 'con khi tu Illinois' --- vao chuc Bo truong Chien tranh.\\
\textit{(He appointed Edwin Stanton --- who had publicly mocked him as 'the gorilla from Illinois' --- as Secretary of War.)}

Cac cong su kinh ngac hoi tai sao; Lincoln tra loi that gian di: 'Ong ta la nguoi gioi nhat cho vi tri nay.'\\
\textit{(His aides were stunned and asked why; Lincoln replied simply: 'He is the best man for the position.')}

Stanton, duoc trong nha nhu vay, tro thanh Bo truong tan tuy nhat trong lich su nuoc My.\\
\textit{(Stanton, treated with such respect, became the most dedicated Secretary in American history.)}

Khi Lincoln bi am sat, chinh Stanton la nguoi khoc va noi: 'Day la nguoi vi dai nhat toi tung biet.'\\
\textit{(When Lincoln was assassinated, it was Stanton who wept and said: 'He is the greatest man I have ever known.')}

\end{truyen}
\begin{baihoc}
  Tha thu ke thu va trao ho co hoi la cach bien ke thu thanh dong minh --- va viet nen lich su.\\
  \textit{(Forgiving an enemy and giving them a chance is how you transform foes into allies --- and make history.)}
\end{baihoc}

\section{Mandela Ra Tu Khong Cam Thu}
\begin{truyen}{Mandela Ra Tu Khong Cam Thu}{Nelson Mandela --- Nam Phi 1990}
\chuhoa{S}{}S au 27 nam bi giam can trong nha tu Robben Island vi dau tranh chong phan biet chung toc, Nelson Mandela buoc ra tu.\\
\textit{(After 27 years imprisoned in Robben Island for fighting apartheid, Nelson Mandela walked free.)}

Phong vien hoi ong co cam thu nhung ke da giam cam ong suot gan ba muoi nam khong.\\
\textit{(A reporter asked if he hated those who had imprisoned him for nearly thirty years.)}

Mandela tra loi: 'Neu toi buoc ra kia voi long cam thu, toi se van con la tu nhan, du ma cua tu da mo tra.'\\
\textit{(Mandela answered: 'If I walked out there with hatred, I would still be a prisoner, even if the prison door had opened.')}

Ong khong chi tha thu --- ong moi nguoi canh sat tung giam giu ong den du le nhau chuc tong thong.\\
\textit{(He didn't just forgive --- he invited the guard who had once imprisoned him to attend his presidential inauguration.)}

Nam Phi thoat khoi vung xoay bao luc nho chinh su khoan dung phi thuong do cua Mandela.\\
\textit{(South Africa escaped a cycle of violence largely because of Mandela's extraordinary forgiveness.)}

\end{truyen}
\begin{baihoc}
  Long cam thu chi giam cam ban trong bung; su tha thu giai phong chinh ban truoc khi giai phong ke thu.\\
  \textit{(Hatred imprisons you inside yourself; forgiveness liberates you before it liberates your enemy.)}
\end{baihoc}

\section{Khong Tu Tha Hoc Tro Phan Boi}
\begin{truyen}{Khong Tu Tha Hoc Tro Phan Boi}{Khong Tu --- Trung Quoc TK 5 TCN}
\chuhoa{M}{}M ot hoc tro cua Khong Tu da cong khai bat vua dung loi khuyen cua su phu de kiem loi ca nhan.\\
\textit{(A student of Confucius openly urged a lord to use the master's advice for personal gain.)}

Cac mon sinh khac yeu cau Khong Tu duoi hoc tro nay di; nha trien hoc im lang hoi tham.\\
\textit{(Other disciples demanded Confucius expel the student; the philosopher calmly asked questions first.)}

Nghe cau tra loi, Khong Tu noi: 'Nguoi ay chua hieu dao; day la loi ta, khong phai loi nguoi ay.'\\
\textit{(After hearing the answer, Confucius said: 'He has not understood the Way yet; that is my fault, not his.')}

Ong tiep tuc day hoc tro do them ba nam cho den khi nguoi nay that su hieu va thay doi.\\
\textit{(He continued teaching that student for three more years until the person truly understood and changed.)}

Hoc tro nay ve sau tro thanh mot trong nhung quan chuc cong bang nhat cua thoi dai do.\\
\textit{(That student later became one of the most just officials of that era.)}

\end{truyen}
\begin{baihoc}
  Nguoi phan boi ban thuong la nguoi chu truoc day ban da khong day du; hay nhan trach nhiem truoc khi phan xet.\\
  \textit{(Someone who betrays you is often someone you didn't teach enough before; take responsibility before judging.)}
\end{baihoc}

\section{Vua Ly Thai To Tha Ke Thu Cu}
\begin{truyen}{Vua Ly Thai To Tha Ke Thu Cu}{Ly Thai To --- Viet Nam 1010}
\chuhoa{K}{}K hi Ly Cong Uan len ngoi thanh lap vuong trieu Ly, nhieu ke tung chong doi ong van con song.\\
\textit{(When Ly Cong Uan ascended the throne and founded the Ly Dynasty, many who had opposed him still lived.)}

Cac quan de nghi truy sat de tru hau hoa; nha vua im lang truoc roi quay sang van de khac.\\
\textit{(Officials proposed hunting them down to prevent future threats; the king fell silent then turned to other matters.)}

Cuoi nam do, ong ban chinh sach an xa rong rai --- tha tat ca toi pham chinh tri va moi nguoi tro ve.\\
\textit{(That year's end, he announced a broad amnesty --- pardoning all political offenders and welcoming all to return.)}

Mot vi quan co cu tra lai va phuc vu Ly Thai To trung thanh den cuoi doi, lam nen nhieu cong tich.\\
\textit{(A former official returned and served Ly Thai To faithfully for life, accomplishing many great things.)}

Trieu dai Ly ke sang mot trong nhung giai doan hung cuong nhat lich su Viet Nam.\\
\textit{(The Ly Dynasty went on to become one of the most prosperous periods in Vietnamese history.)}

\end{truyen}
\begin{baihoc}
  Bau du khoan dung bao gio cung tao ra mot dat nuoc manh hon la bau du nuoc mat va thu han sinh ra.\\
  \textit{(Sowing forgiveness always creates a stronger nation than sowing resentment and hatred.)}
\end{baihoc}

\section{Corrie Ten Boom Tha Ten Linh Duc}
\begin{truyen}{Corrie Ten Boom Tha Ten Linh Duc}{Corrie ten Boom --- Duc 1947}
\chuhoa{S}{}S au Chien tranh The gioi II, Corrie ten Boom--- nguoi phu nu Ha Lan da song sot sau trai tap trung --- den Duc de giang ve su tha thu.\\
\textit{(After World War II, Corrie ten Boom --- a Dutch woman who survived the concentration camps --- came to Germany to speak about forgiveness.)}

Sau bai giang, mot nguoi dan ong buoc den bat tay ba; Corrie nhan ra day chinh la ten linh da bao hanh ba trong trai.\\
\textit{(After the lecture, a man approached to shake her hand; Corrie recognized him as the very guard who had brutalized her in the camp.)}

Tim ba dong cung, tay ba khong muon dua ra, nhung ba van cu cau nguyen va dua tay biet ra.\\
\textit{(Her heart froze and her hand refused to move, but she prayed and slowly extended her hand anyway.)}

Ngay khi tay ba cham vao tay ong, ba cam nhan su binh an tuon chay qua nguoi nhu mot dong dien.\\
\textit{(The moment her hand touched his, she felt peace flowing through her like an electric current.)}

Ba viet lai: 'Tinh yeu la thu manh me hon bat ky cam xuc nao khac --- va su tha thu la co the chon lua.'\\
\textit{(She wrote: 'Love is stronger than any feeling --- and forgiveness is a choice that can be made.')}

\end{truyen}
\begin{baihoc}
  Su tha thu chua lan luot ca nguoi tha thu lan nguoi duoc tha; no an me boi vung tuong gan nang ne nhat.\\
  \textit{(Forgiveness heals both the one who forgives and the one forgiven; it opens the heaviest gates.)}
\end{baihoc}

\section{Gandhi Va Nguoi Linh Bao Hanh}
\begin{truyen}{Gandhi Va Nguoi Linh Bao Hanh}{Mahatma Gandhi --- An Do 1922}
\chuhoa{N}{}N am 1922, Gandhi bi bat giam va bi dua ra toa vi toi xui gio cuoc noi loan chong thuc dan Anh.\\
\textit{(In 1922, Gandhi was arrested and brought to trial for the charge of inciting rebellion against British rule.)}

Thay quan toa - mot nguoi Anh ten Broomfield --- nói rang ong rat kho xu voi truong hop nay.\\
\textit{(The judge --- an Englishman named Broomfield --- said he found this case deeply difficult to handle.)}

Broomfield noi: 'Neu phap luat cho toi, toi se tha ong di; nhung luat phap bat toi phai xu ong.'\\
\textit{(Broomfield said: 'If the law allowed me, I would set you free; but the law requires me to sentence you.')}

Gandhi mim cuoi dap: 'Thua quy toa, toi biet quy toa da lam tat ca nhung gi quy toa co the lam.'\\
\textit{(Gandhi smiled and replied: 'Your Honor, I know you have done everything within your power.')}

Chinh su khoan dung lan nhin nhan su khoan dung do da lam mem long nhieu nguoi Anh hon bat ky cuoc bieu tinh nao.\\
\textit{(That mutual recognition of respect moved more English hearts than any protest march ever had.)}

\end{truyen}
\begin{baihoc}
  Khoan dung voi ke đang bat minh la hanh dong manh me nhat co the thuc hien truoc mat mot the gioi dang nhin.\\
  \textit{(Being gracious to those who oppress you is the most powerful act you can perform before a watching world.)}
\end{baihoc}

\section{Tong Tu Truong --- Chi Mot Noi Lan}
\begin{truyen}{Tong Tu Truong --- Chi Mot Noi Lan}{Truyen Co Trung Hoa --- TK 6 TCN}
\chuhoa{T}{}T ong Tu Truong la quan dai phu nuoc te thoi Xuan Thu co tieng la nguoi hieu lua, biet nguoi.\\
\textit{(Tong Ziqi was a high official of the state of Qi during the Spring and Autumn period, famed for understanding people.)}

Mot nguoi quan chuc cap duoi da vu cao ong truoc vua de gianh chuc vi cua ong.\\
\textit{(A subordinate official had falsely accused him before the king to seize his position.)}

Sau khi sac thu chung minh ong vo toi, moi nguoi doi biet ong se tra thu ke vu cao.\\
\textit{(When evidence proved his innocence, everyone waited to see him take revenge on his accuser.)}

Ong khong lam vay; thay vao do ong chu dong de nghi vua khoan hong cho ke do, noi rang: 'Ke do chi lam vi qua khat vong.'\\
\textit{(He did not; instead he proactively asked the king to pardon the man, saying: 'He only did it out of too much ambition.')}

Nguoi quan chuc nay sau do tro thanh nguoi ung ho nhiet tinh nhat cua ong, biet on suot ca doi.\\
\textit{(The official afterward became his most zealous supporter, grateful for the rest of his life.)}

\end{truyen}
\begin{baihoc}
  Kho khoan dung hon sang tac; nhung phan thuong cua no -- long trung thanh -- khong tien nao mua duoc.\\
  \textit{(Forgiveness is harder than revenge; but its reward --- loyalty --- cannot be bought with any price.)}
\end{baihoc}

\section{Archbishop Tutu Va Uy Ban Hoa Giai}
\begin{truyen}{Archbishop Tutu Va Uy Ban Hoa Giai}{Desmond Tutu --- Nam Phi 1995}
\chuhoa{S}{}S au khi che do apartheid sup do, Nam Phi dung truoc lua chon: xet xu toan bo hay hoa giai.\\
\textit{(After apartheid collapsed, South Africa faced a choice: total prosecution or reconciliation.)}

Archbishop Desmond Tutu dung ra lanh dao Uy Ban Su That va Hoa Giai --- mot thi nghiem chua tung co.\\
\textit{(Archbishop Desmond Tutu led the Truth and Reconciliation Commission --- an unprecedented experiment.)}

Nhung ke hanh hung phai ra khai nhan toi truoc cong chung; nan nhan duoc quyen noi ve noi dau cua ho.\\
\textit{(Perpetrators had to confess their crimes publicly; victims were given the right to speak of their pain.)}

Nhieu nan nhan --- nhung nguoi co quyen cam thu nhat --- da chon su tha thu thay vi tra thu.\\
\textit{(Many victims --- those with the greatest right to hatred --- chose forgiveness over revenge.)}

Uy Ban nay giup Nam Phi tranh duoc mot cuoc noi chien ma nhieu nguoi du doan se khong the tranh khoi.\\
\textit{(The Commission helped South Africa avoid a civil war that many had predicted was unavoidable.)}

\end{truyen}
\begin{baihoc}
  Su that khong co tha thu se tao nen bao luc; tha thu khong co su that se la lua doi --- chi co ca hai moi chua lanh.\\
  \textit{(Truth without forgiveness creates violence; forgiveness without truth is deception --- only both together can heal.)}
\end{baihoc}

\section{Hoang De Minh --- Dai An Xa Cho Sinh Linh}
\begin{truyen}{Hoang De Minh --- Dai An Xa Cho Sinh Linh}{Minh Thai To --- Trung Hoa 1368}
\chuhoa{S}{}S au khi thanh lap trieu Minh va danh duoi nguoi Mong Co, Chu Nguyen Chuong dung truoc yeu cau cua quan lanh.\\
\textit{(After founding the Ming Dynasty and expelling the Mongols, Zhu Yuanzhang faced demands from his generals.)}

Ho thuc ep ong tru diet toan bo gia dinh cac tong binh Mong Co de tru hau hoa mai sau.\\
\textit{(They pressed him to exterminate the entire families of Mongol commanders to prevent future threats.)}

Hoang de hoi cac binh linh Mong Co con lai: 'Cac nguoi se trung thanh neu nhan duoc khoan dung?'\\
\textit{(The emperor asked the remaining Mongol soldiers: 'Will you be loyal if you receive mercy?')}

Khi ho ban hon, ong ra lenh bo vu khi di va trao dat de canh tac --- khong giet them mot ai nua.\\
\textit{(When they wavered, he ordered them to lay down arms and gave them land to farm --- killing no one more.)}

Bien gioi phia bac giu yen tinh trong nhieu thap nien nhot nho chinh sach an hoa bat ngo nay.\\
\textit{(The northern border remained peaceful for decades precisely because of this unexpected policy of mercy.)}

\end{truyen}
\begin{baihoc}
  Mot quy dao tha thu co the mang lai hoa binh lau dai hon mot tram tran chien thang quan su.\\
  \textit{(One royal act of forgiveness can bring longer peace than a hundred military victories.)}
\end{baihoc}

\section{Vo Nguyen Giap Tha Tu Binh Phap}
\begin{truyen}{Vo Nguyen Giap Tha Tu Binh Phap}{Vo Nguyen Giap --- Viet Nam 1954}
\chuhoa{S}{}S au chien thang Dien Bien Phu nam 1954, Viet Nam co hang ngan tu binh Phap trong tay.\\
\textit{(After the victory at Dien Bien Phu in 1954, Vietnam held thousands of French prisoners of war.)}

Tuong Vo Nguyen Giap ra lenh doi xu nhan dao voi tu binh --- cham soc thuong binh, cung cap thuc an du.\\
\textit{(General Vo Nguyen Giap ordered humane treatment for prisoners --- caring for the wounded, providing sufficient food.)}

Nhieu van phong nguoc chieu mong doi su bao thu; nhung de nghi tra thu bi tuong tu choi.\\
\textit{(Many historical accounts note expectations of retaliation; but requests for revenge were refused by the general.)}

Mot so si quan Phap ve nuoc da viet sach ke lai su doi xu co nhan nay voi long kinh trong sau sac.\\
\textit{(Some French officers returned home and wrote books recounting this humane treatment with deep respect.)}

Su khoan hoa nay gop phan xay dung hinh anh dau tranh chinh nghia cua Viet Nam truoc mat the gioi.\\
\textit{(This magnanimity helped build Vietnam's image as a righteous struggle in the eyes of the world.)}

\end{truyen}
\begin{baihoc}
  Doi xu nhan dao voi ke thua tran khong chi la dao duc --- do con la chien luoc ngoai giao khon ngoan nhat.\\
  \textit{(Treating defeated enemies humanely is not only moral --- it is the wisest diplomatic strategy.)}
\end{baihoc}
